\section{Tilted array predictions and unordered tuple exceptions}\label{tra:tilted-arrays}

We give a survival-tilted proof of the weighted compatibility interface
from Section~\ref{tra:duality}. It retains the integrability costs of
arbitrary line weights through marginal deficits and entropy duality.
The comparison distribution is tilted by reciprocal survival
probabilities, giving another way to charge the exceptional atoms.

At each cell the proof multiplies the probabilities of available light symbols by the reciprocal survival probability for symbol pairs located in other columns. Pairs in the current column are exceptions, whose contribution is included in the normalizing cost. The logarithm of the multiplier supplies the entropy gain. Its normalizing cost remains small after fixing the full compatible array, so rare heavy atoms are charged by their number without an additional logarithmic loss.

The sparse step retains a second piece of information. An average interaction estimate identifies exceptional input tuples, but an operator norm also requires control of how much mass those tuples can receive from every output tuple. We prove that control for every exceptional set determined by unordered input positions, then use it in two Schur bounds.

Here $N=2^d$, $d\ge1$. Write $T_N$ for a sweep and $\mathcal Z_N(u)=\operatorname{Tr}(T_N^*T_N)^u$ for $u>0$, with unnormalized regular trace. One sweep costs $d$ physical shuffles as in Section~\ref{sec:prelim}. Products act right-to-left, as in Section~\ref{sec:prelim}; a row move followed by a column move therefore places the column operator on the left. All coefficient functions below are densities relative to uniform probability on the indicated subgroup.

\subsection{Weighted compatibility of two permutation arrays}
\label{tra:tilted-arrays:array-definition}
We first prove the weighted array inequality independently of the sweep. Consider an \(n\)-by-\(m\) grid with \(m\le n\le 2m\); in this estimate write \(N=nm\). Arrays \(A,B\) satisfy:
\begin{itemize}
\item each row \(A_i\) of \(A\) gives the \(m\) symbols once each;
\item each column \(B_j\) of \(B\) gives the \(n\) symbols (of a second alphabet) once each.
\end{itemize}
\begin{proposition}\label{tra:tilted-arrays:weighted}
\label{tra:tilted-arrays:weighted-statement}
Use as reference \(U\) independent uniform permutations in these rows and columns. Say \(A,B\) are \textbf{compatible} if the symbol pairs \((A_{ij},B_{ij})\) are all distinct. There are absolute \(C_0,C_1,\varepsilon>0\) such that, for \(m\) sufficiently large, putting
\[
\alpha=1-\frac{C_0}{\log m}>0,
\]
we have the following bound on nonnegative weights, with uniform expectations on its right:
\begin{equation}
\begin{split}
\mathbf E_U\left[\mathbf 1_{\mathrm{comp}}\prod_i L_i(A_i)\prod_j M_j(B_j)\right]
\ \le\ &\exp\{-N+C_1N m^{-\varepsilon}\}\\
&{}\cdot
\prod_i\left(\mathbf E L_i^{1/\alpha}\right)^\alpha
\prod_j\left(\mathbf E M_j^{1/\alpha}\right)^\alpha .
\end{split}\label{tra:tilted-arrays:eq2}
\end{equation}
\end{proposition}
\begin{proof}\label{tra:tilted-arrays:weighted-proof}
\label{tra:tilted-arrays:tilted-entropy-proof}
To prove this, let \(q\) be any probability law on compatible arrays, with the entropy \(H\) below computed in this law. Define the deficits
\[
D_A=n\log(m!)-H(A),
\qquad
D_B=m\log(n!)-\sum_j H(B_j).
\]
In particular \(D_A\) is for the full \(A\) array (so dominates the sum of marginal deficits in the rows). Let
\[
I=\sum_j H(B_j)-H(B\mid A).
\]
The column marginal entropies here are unconditional. Thus $I$ includes
both the conditional correlation of the columns given $A$ and the sum
of their separate mutual informations with $A$, as in the comparison
after Lemma~\ref{lem:conditional-column-entropy}.
We claim
\begin{equation}
 I\ \ge\ N-C_1N m^{-\varepsilon}
             -\frac{C_0}{\log m}(D_A+D_B). \label{tra:tilted-arrays:eq3}
\end{equation}

Reveal \(A\) and go through the columns \(j\) of \(B\) in order of iid priorities \(x_j\) uniform on \([0,1]\) (from small to large), these priorities independent of \(A,B\). Within each column use row order \(i=1,\dots,n\). At cell \(i,j\), use the marginal conditional probabilities
\(p_b\) for the column \(B_j\) itself, given its prefix, as comparison. Then \(I\) is the sum over cells of the expected KL divergences of the true conditional (including \(A\), priorities and earlier revealed \(B\)) against \(p\). This follows from the conditional entropy chain rule in this order, since log probabilities from the \(p\)'s are just marginal-column log probabilities in total.

Use the following tilt of \(p\). Say an atom \(b\) here is light if \(p_b\le m^{-3/10}\). Write \(a=A_{ij}\) and let the multiplicity \(v\) be the number of times \(a\) occurs in the \(j\)-th column of \(A\). Write \(y=1-x_j\). If \(v>m^{1/10}\) or \(y<m^{-1/10}\), skip the tilt, using factors \(t_b=1\). Otherwise let \(t_b=1\) on the non-light atoms, and on the light atoms set
\[
 t_b=
 \begin{cases}
  1/y, & \text{pair }(a,b)\text{ not present in any earlier column},\\
  0, &\text{otherwise}.
 \end{cases}
\]
These factors use only data known at this point and the normalizer \(Z=\sum_b p_b t_b\) is positive on histories of positive conditional probability (with comparisons understood for almost every priority choice), by compatibility. Thus the divergence contribution is bounded below in expectation by the expected
\[
 \log t_{B_{ij}}-\log Z.
\]

We can bound the normalizer cost by holding the \emph{full} compatible arrays and \(x_j\) fixed and taking expectation in other priorities. In particular \(p\) has not depended on those priorities, because within-column revelation is in fixed order. Except for the \(v\) symbols paired with \(a\) in column \(j\), any \(b\) has its pair with \(a\) in a different column, so the pair is available (not in an earlier column) with probability \(y\). Consequently in non-skipped cases
\[
 \mathbf E_{\text{other priorities}} Z
 \le 1+v m^{-3/10}/y
 \le 1+m^{-1/10}.
\]
Log-normalizer cost in expectation is then at most \(m^{-1/10}\). On the other hand if multiplicity is not too high and the actual atom is light, we earn in expectation over priorities
\(\int_{m^{-1/10}}^1 \log(1/y)\,dy
=1-O(m^{-1/10}\log m)\).
This use of column priorities lets the potential loss from skipping non-light atoms be counted without a large logarithmic penalty per lost cell.

The tilt has supplied almost one nat at each retained light cell, while its expected normalizing cost is small. Two exclusions remain: a heavy actual symbol, and a symbol occurring too often in its column of $A$. We pay for these using, respectively, $D_B$ and the full-array deficit $D_A$.

\begin{itemize}
\item Non-light actual atoms have expected count
\end{itemize}
\begin{equation}
 O\left(N m^{-c}+\frac{D_B}{\log m}\right) \label{tra:tilted-arrays:eq4}
\end{equation}
for an absolute \(c>0\). Indeed, we may ignore the last \(\lceil m^{4/5}\rceil\) entries in each column, counting those as losses. At any remaining cell, the non-light atoms are at most fraction \(O(m^{-1/2})\) of the alphabet remaining in a uniform permutation. If their conditional \(p\)-mass \(h\) is at least \(m^{-1/4}\), KL divergence versus that uniform choice is \(\Omega(h\log m)\), by divergence on the related binary event (\(h\log(h/p_{\mathrm{unif}}(\mathrm{event}))-h\) suffices for a lower bound). Else the mass is already small. The expected divergences of \(p\) versus the uniform remaining choices sum to \(D_B\), proving \eqref{tra:tilted-arrays:eq4}.

\begin{itemize}
\item For \(A\), let \(J\) count cells at which the multiplicity in the column of the cell's symbol exceeds \(m^{1/10}\). We have
\end{itemize}
\begin{equation}
 \mathbf E_q J\le C\left(Nm^{-1/50}+\frac{D_A}{\log m}\right). \label{tra:tilted-arrays:eq5}
\end{equation}
To see it, under independent uniform rows, for integer \(z\ge z_0=\lceil N m^{-1/50}\rceil\),
\[
 \mathbb P(J\ge z)\le \exp(-0.05 z\log m)
\]
for large \(m\). In fact, to witness the event, in each of the \(m\) columns choose the list of symbols having high multiplicity (size at most \(2m^{9/10}\)) and prescribe \(z\) cells to take symbols from the lists of their columns. The log number of choices of lists is \(O(Nm^{-1/10}\log m)\), and cells have log cost at most \(z\log(eN/z)\le z(0.02\log m+1)\). For a row with \(u\) prescribed cells the uniform success probability is bounded by
\((2m^{9/10})^u/(m)_u\le (2e m^{-1/10})^u\), where subscript denotes falling factorial. (We use \((m)_u\ge (m/e)^u\), following e.g. from the factorial bound and prefix geometric means.) These estimates give the probability bound. Applying exponential-moment comparison by relative entropy with parameter \(0.02\log m\) now gives \eqref{tra:tilted-arrays:eq5}: log of the uniform exponential moment is at most \(0.02 z_0\log m+O(1)\), whereas \(D_A\) bounds parameter times expected \(J\) under \(q\) minus this log.

This proves \eqref{tra:tilted-arrays:eq3}, adjusting constants and using, say, a sufficiently small \(\varepsilon\in(0,1/50)\). The full KL divergence of \(q\) relative to \(U\) is \(D_A+D_B+I\), so \eqref{tra:tilted-arrays:eq3} implies \eqref{tra:tilted-arrays:eq2} by entropy variational comparison: the divergence pays \(N-C_1 N m^{-\varepsilon}\) and still pays \(\alpha\) times the sum of the row and column marginal deficits, against which expected log weights are bounded using \(\alpha\log\mathbf E L_i^{1/\alpha}\), \(\alpha\log\mathbf E M_j^{1/\alpha}\). More explicitly, for a weight \(L_i\),
\(\mathbf E_q\log L_i\) minus \(\alpha\) times marginal deficit is at most the former of these log expressions, and likewise for \(M_j\); then take the supremum over \(q\) for the left side log of \eqref{tra:tilted-arrays:eq2}. These variational facts follow as usual by nonnegativity of divergence from a tilted law; weights with zeros are allowed by limits or restriction. This completes the grid estimate.
\end{proof}
\subsection{The intermediate singular moment}
\label{tra:tilted-arrays:operator-split}
The weighted inequality applies to squared absolute coefficient densities. We now check their norm exponents, so that the child moment bounds produce a weak parent moment with only a factor $1+O(1/d)$ increase in the exponent.

Take \(m=2^{\lfloor d/2\rfloor}\), \(n=2^{\lceil d/2\rceil}\), for \(d\) large, and assign positions to an \(n\)-by-\(m\) grid so the bits for the column are the first \(\lfloor d/2\rfloor\) bits. Write \(\mathcal L\) for the group of row-preserving permutations, \(\mathcal R\) for the column-preserving group. On the regular representation,
\[
 T_N=K_{\mathcal R}K_{\mathcal L},
\]
where $K_{\mathcal L}$ uses independent $T_m$ sweeps on the rows and $K_{\mathcal R}$ uses independent $T_n$ sweeps on the columns. Set $X=K_{\mathcal L}K_{\mathcal L}^*$ and $Y=K_{\mathcal R}^*K_{\mathcal R}$. Thus $T_N^*T_N=K_{\mathcal L}^*YK_{\mathcal L}$ has the same eigenvalues, including zeros, as $Y^{1/2}XY^{1/2}$.
\begin{proposition}\label{tra:tilted-arrays:recursion}
\label{tra:tilted-arrays:recursion-statement}
Suppose
\[
 \mathcal Z_m(u_m)\le 6/5,\qquad \mathcal Z_n(u_n)\le 6/5
\]
with \(u_m,u_n\ge 2\), and, using \(\alpha\) from \eqref{tra:tilted-arrays:eq2}, set
\[
 \bar r=(2-\alpha)\max(u_m,u_n),\qquad p=\bar r/2 .
\]
Then for some absolute \(\eta\in(0,1)\),
\begin{equation}
 \mathcal Z_N(\bar r)\le \exp\left(N^{1-\eta}\right). \label{tra:tilted-arrays:eq6}
\end{equation}
\end{proposition}
\begin{proof}\label{tra:tilted-arrays:recursion-proof}
\label{tra:tilted-arrays:fourth-trace-and-inverse-transform}
Lemma~\ref{lem:positive-power-comparison}, with $\bar r=2p\ge2$,
gives
\begin{equation}
 \mathcal Z_N(\bar r)
 \le\operatorname{Tr}(X^pY^pX^pY^p).
 \label{tra:tilted-arrays:eq7}
\end{equation}

Within their respective subgroups express \(X^p,Y^p\) by densities \(F,G_1\) on the uniform probability measures (a density here is just a coefficient function, not required to be nonnegative). They are products over rows, respectively columns, of the individual densities for the smaller problems: \(F(l)\) uses a factor \(f(l_i)\) for each of the row permutations \(l_i\) of \(l\in\mathcal L\), and \(G_1(r)\) uses a factor \(f_1(r_j)\) for each column permutation of \(r\in\mathcal R\). Here $f$ is for $(T_mT_m^*)^p$, and $f_1$ for $(T_n^*T_n)^p$. Both orientations have the child trace moments appearing in the hypothesis. This follows by tensor products and the fact that restriction of the regular action to a subgroup gives copies of its regular action, so these powers in a subgroup convolution algebra simply embed in the full group even when using real positive powers.

Apply Lemma~\ref{lem:inverse-fourier-bound} to each individual
density with $a=2/\alpha$ and $a'=2/(2-\alpha)$. For the row density,
the regular Fourier moment on its right is
$\mathcal Z_m(pa')$, with
\[
 pa'=\bar r/(2-\alpha)=\max(u_m,u_n)\ge u_m.
\]
The column calculation uses $\mathcal Z_n(pa')$ and $pa'\ge u_n$.
Since the sweep operators are contractions, we get
\begin{equation}
 \|f\|_{2/\alpha},\ \|f_1\|_{2/\alpha}\ \le 2. \label{tra:tilted-arrays:eq8}
\end{equation}

Use Lemma~\ref{lem:coefficient-compatibility}, in its unnormalized
form \eqref{e:four-factor}, on the positive line factors of $X^p,Y^p$.
It gives
\begin{equation}
 \mathcal Z_N(\bar r)
 \le\frac{|S_N|}{|\mathcal L||\mathcal R|}
 \mathbb E_{l,r}\!\left[
   \mathbf1_{rl\in\mathcal L\mathcal R}|F(l)|^2|G_1(r)|^2\right],
 \label{tra:tilted-arrays:eq9}
\end{equation}
where $l,r$ are independent subgroup uniforms. In the grid between
the row move $l$ and the column move $r$, the original-column labels
are $A_{ij}=l_i^{-1}(j)$ and the destination-row labels are
$B_{ij}=r_j(i)$. The event in \eqref{tra:tilted-arrays:eq9} is exactly
that all pairs $(A_{ij},B_{ij})$ are distinct. The row weights expressed
in these arrays are $|f(A_i^{-1})|^2$, whose uniform norms equal those
of $|f|^2$; the column weights are $|f_1(B_j)|^2$.
Thus the weighted estimate \eqref{tra:tilted-arrays:eq2} and
\eqref{tra:tilted-arrays:eq8} give
\[
 \log\mathbb E_{l,r}\!\left[
   \mathbf1_{rl\in\mathcal L\mathcal R}|F(l)|^2|G_1(r)|^2\right]
 \le-N+C_1Nm^{-\varepsilon}+C(n+m).
\]

Also
\(\log[ |S_N|/(|\mathcal L||\mathcal R|)]\le N+O(\log N)\)
by applying the upper bound $\log(N!)\le N\log N-N+O(\log N)$ to the numerator and the lower bounds $\log(m!)\ge m\log m-m$, $\log(n!)\ge n\log n-n$ to the denominator. This one-sided bound is all the recursion needs. This proves \eqref{tra:tilted-arrays:eq6}, taking a sufficiently small fixed positive \(\eta\) (e.g. less than \(\varepsilon/3\) and \(1/4\)), for all sufficiently large \(d\).
\end{proof}
\subsection{Unordered bad rows at sparse levels}
The weak regular moment controls representations of large dimension. The remaining input is an operator bound when the first row is long. We first cancel isolated paths, then estimate the mean of the resulting absolute kernel, and finally control exceptional input rows in every output column.
\begin{lemma}\label{tra:tilted-arrays:sparse}
\label{tra:tilted-arrays:sparse-statement}
Fix \(0<\delta<1\). We need the following bound: if \(\lambda_1=N-k\) where
\(1\le k\le N^{1-\delta}\), then the operator norm in this representation of \(T_N\) is at most
\begin{equation}
 N^{-b k} \label{tra:tilted-arrays:eq10}
\end{equation}
for some \(b=b(\delta)>0\) and \(d\) sufficiently large. On the sign twist it is even zero (this latter statement also holds for \(k=0\)).
\end{lemma}
\begin{proof}\label{tra:tilted-arrays:sparse-proof}
\label{tra:tilted-arrays:forest-and-unordered-set-proof}
To prove this, let \(\mathcal X\) be ordered \(k\)-tuples of distinct positions, measure \(\nu\) uniform on \(\mathcal X\). Functions on \(\mathcal X\) are acted on by position permutations. The shape \(\lambda\) occurs here and does not occur in functions of any proper subset of the coordinates. Indeed occurrence on \(k\)-tuples is having invariants under \(S_{N-k}\), by the action on tuples and its stabilizer, and the branching rule says this holds for shapes with first row at least \(N-k\). Smaller coordinate subsets would require a larger first row.

For \(S\subseteq [k]\), put \(s=|S|\), and for \(x,y\in\mathcal X\) let
\[
 R_S(x,y)=N^s\,\mathbb P(g(x_S)=y_S),
\]
using the pass permutation. The full \(R_{[k]}\) as an integral kernel against \(\nu\) represents \(T_N^*\) on tuple functions (the forward transition kernel is the adjoint under our action convention), up to the factor \(N^k/(N)_k\). Also, in bilinear forms on a copy of \(\lambda\), the proper subset terms give zero since they only involve those subsets of coordinates. Thus we can bound the desired norm by giving an operator norm bound on
\[
 Z_*(x,y)=\sum_{S\subseteq [k]}(-1)^{k-|S|}R_S(x,y) .
\]

A path of a card from one address to another in the \(d\)-layer pass is unique: at each coordinate it must go to the bit given by its destination address. It occurs with probability \(1/N\); it is prescribed by individual switch choices. Make a graph on \([k]\) using \(x,y\), putting an edge whenever the two associated paths use the same switch at some layer (in the same or different directions, including use of the same switch input). If there is an isolated vertex, \(Z_*=0\). In fact such a path then imposes switch choices independent of those for any other subset of paths, which gives cancellation in pairs due to the scaling by \(N\). Writing \(\Gamma\) for the event the graph has no isolated vertex, define
\[
 W(x)=\int\sum_S R_S(x,y)\mathbf 1_\Gamma\,d\nu(y),
\]
an upper bound on the absolute row integral. We have \(W(x)\le e^{Ck}\), since the row integral of each \(R_S\) is \(N^s/(N)_s\le e^s\). Likewise \(e^{Ck}\) bounds columns using the inverse pass. We will obtain much smaller bounds by controlling typical rows and showing their exceptions cannot concentrate in any column.

First, for some \(\gamma=\gamma(\delta)>0\),
\begin{equation}
 \mathbf E_\nu W\le N^{-\gamma k}. \label{tra:tilted-arrays:eq11}
\end{equation}
For \(k=1\) the no-isolated-vertex event is impossible, so take \(k\ge 2\). In bounding each average over \(x,y\), we can take all \(x_i,y_i\) iid addresses and relax distinctness at cost at most \(e^{2k}\), using the same \(R_S\) definition also on repeats. Weighting that iid law by \(R_S\) gives a simple sampling rule: sample a full pass network, take all \(x_i\) iid, use \(y_i=g(x_i)\) on \(S\), and sample the other \(y_i\) iid. Conditional on this network the resulting single-card paths are independent. At each layer the switch used by each such path is uniform among \(N/2\) switches: on \(S\) this is by uniform start through the fixed network, and outside \(S\) by iid start and finish (the address along the path uses respective bits of these). Therefore for any forest on the cards with \(u\) edges, the probability its edges are all path interactions is at most \((2d/N)^u\): for a leaf relative to its parent path this bound per edge follows by union over layers, and leaves can be peeled off using independence conditional on network. If \(\Gamma\) holds, there is a forest with \(u\) between \(k/2\) and \(k-1\) edges. At each such size there are at most \((C k)^u\) to try, by choosing edges. Hence we have upper bound
\[
 2^k e^{2k}\sum_{\lceil k/2\rceil\le u\le k-1}(C k d/N)^u
\]
for \eqref{tra:tilted-arrays:eq11}, proving the claim for the indicated range of \(k\).

The mean estimate alone does not control a matrix norm: a small collection of input rows might carry a large operator. The next domination estimate rules this out in each fixed output column.

Next, for any set \(\mathcal B\) of \(x\)'s depending only on the unordered set of their entries, we have, for every fixed \(y\in\mathcal X\),
\begin{equation}
 \int_{\mathcal B} R_S(x,y)\,d\nu(x)\le e^{Ck}\nu(\mathcal B). \label{tra:tilted-arrays:eq12}
\end{equation}
To see it we can again use iid \(x\) at cost at most \(e^k\), asking also for distinctness along with \(\mathcal B\). Weighting by \(R_S\) amounts here to setting \(x_S=g^{-1}(y_S)\) and sampling remaining \(x_i\) iid. Under this inverse network, the set occupied by the \(s\) paths satisfies, for each set of sites \(T\) at any stage, probability of containing \(T\) at most product over \(T\) of the current site marginals. Indeed this property holds initially for the deterministic set at \(y_S\), and remains true through each fair random transposition, which suffices as these give the layers. For \(T\) using one of the switched sites, average the previous bounds with that site exchanged, and for using both, use that product of the two site marginals is bounded by product after averaging the two; other marginals have not changed. At the end the site marginals are \(s/N\) by the uniform single-position transition. Thus probability for the set of \(s\) sites to equal any particular \(s\)-set is at most \((s/N)^s\), interpreted as 1 for \(s=0\). For all of \(x\) to have as entries a given \(k\)-set, distinctly, in this law, probability is thus at most
\[
 \binom{k}{s}(s/N)^s (k-s)!/N^{k-s}
 \le e^k k!/N^k .
\]
The probability for that set under \(\nu\) is \(k!/(N)_k\), giving \eqref{tra:tilted-arrays:eq12} as desired.

Now \(W(x)\) is unordered-set measurable by simultaneous reindexing of \(x,y\), so the rows with \(W>N^{-\gamma k/2}\) make a bad set of the kind in \eqref{tra:tilted-arrays:eq12}, having \(\nu\)-probability at most \(N^{-\gamma k/2}\). For \(Z_*\) restricted to good rows, its absolute row and column integrals are at most \(N^{-\gamma k/2}, e^{Ck}\), respectively; on bad rows, bounds \(e^{Ck},e^{C'k} N^{-\gamma k/2}\) apply respectively by \eqref{tra:tilted-arrays:eq12} and \(|Z_*|\le \sum_S R_S\). Schur's operator norm bound (square root of product of row and column bounds) proves \eqref{tra:tilted-arrays:eq10}, decreasing positive constants in exponents as needed. Here in passing to the copy of \(\lambda\) we indeed multiply the integral kernel by at most 1 to account for the earlier factor \(N^k/(N)_k\).

Finally if we twist the tuple representation by sign, for \(k<N/2\) averaging even one network layer already gives the zero operator. For each tuple there is a switched pair fixing all its positions (i.e. using two unoccupied positions), so the signed subgroup average there cancels. This proves the sign twist claim.
\end{proof}
\subsection{A bounded recursive exponent}
\begin{theorem}\label{tra:tilted-arrays:moment}
\label{tra:tilted-arrays:main-statement}
There is a sequence $r(d)\ge2$ with $\sup_d r(d)<\infty$ such that
\begin{equation}
\mathcal Z_N(r(d))\le 6/5
\qquad(d\ge1). \label{tra:tilted-arrays:eq1}
\end{equation}
For every deterministic initial deck and every integer $w\ge\sup_d r(d)$, the law after $wd$ physical shuffles has chi-square divergence at most $1/5$ from uniform, and hence total-variation distance at most $\frac12\sqrt{1/5}$ from uniform.
\end{theorem}
\begin{proof}[Proof of Theorem~\ref{tra:tilted-arrays:moment}]
\label{tra:tilted-arrays:moment-induction}
The weak moment and the sparse operator bound are independent inputs. We now show that they cover every representation and that the exponent increases at successive splits have a bounded product.

We spell out which representations are covered by \eqref{tra:tilted-arrays:eq10} and its sign version, taking
\(\delta=\eta/4\) for the \(\eta\) in \eqref{tra:tilted-arrays:eq6}. Write \(D_\lambda\) for dimensions. For large \(d\), unless
\begin{equation}
 \log D_\lambda\ \ge\ N^{1-\eta/2}, \label{tra:tilted-arrays:eq13}
\end{equation}
we necessarily have
\begin{equation}
 k=N-\max(\lambda_1,\lambda^{\mathsf t}_1)\ \le\ N^{1-\delta}. \label{tra:tilted-arrays:eq14}
\end{equation}
Here is one way to check. Put \(h=\max(\lambda_1,\lambda^{\mathsf t}_1)\), and if \(h\le N/10\), the hook-length formula already gives at least \((5/e)^N\) from hooks at most \(2h\). Else we can transpose if needed and take \(h\) as the first row length. If there are at least \(z=\lfloor N^{1-\delta}\rfloor\) boxes below the row, take a subdiagram consisting of the row and just \(z\) such boxes. The dimension of the original is at least the number of standard fillings of this subdiagram, which is at least \(\binom{h}{z}\). In fact one can fill the first \(z\) boxes of the top row first (having \(z<h\)), then freely interleave its remainder with the lower part in a fixed standard order on that part, since all lower boxes are in columns at most \(z\); and fillings extend to the original. This suffices for \eqref{tra:tilted-arrays:eq13}, proving \eqref{tra:tilted-arrays:eq14} when \eqref{tra:tilted-arrays:eq13} fails.

Choose a fixed $r_{\min}\ge2$ with $2br_{\min}\ge4$, where $b$
is the sparse constant for $\delta=\eta/4$. At each level $k$ in
\eqref{tra:tilted-arrays:eq14} there are at most $2^{k+1}$ shapes,
allowing transpose, and each has dimension at most $N^k$.
The sign twists are annihilated by the sparse lemma; for the other
nonconstant shapes, \eqref{tra:tilted-arrays:eq10} bounds their regular
contribution at every $u\ge r_{\min}$ by
\[
 \sum_{1\le k\le N^{1-\delta}}
       2^{k+1}N^{2k-2bku}=o(1).
\]
The sign representation itself also contributes zero. These are the
first class in Lemma~\ref{lem:balanced-moment-closure}, with
$\varepsilon=1/5$; their contribution is at most $1/10$ above an
absolute threshold.

The remaining class satisfies \eqref{tra:tilted-arrays:eq13}.
Use the positive-square convention $\nu=2$ and set
\[
 E_N=N^{1-\eta},\qquad H_N=N^{1-\eta/2},\qquad
 \alpha_d=2-\alpha,\qquad\delta_d=d^{-1/2}.
\]
Proposition~\ref{tra:tilted-arrays:recursion} supplies its child-to-parent
hypothesis uniformly for inherited exponents at least $r_{\min}$.
The decisive comparison is
\[
 (1+\delta_d)E_N-\delta_dH_N
 =(1+d^{-1/2})N^{1-\eta}-d^{-1/2}N^{1-\eta/2}
 \longrightarrow-\infty.
\]
Also $(2-\alpha)(1+d^{-1/2})=1+O(d^{-1/2})$.
The common lemma therefore supplies bounded recursive exponents
$r(d)$ satisfying \eqref{tra:tilted-arrays:eq1}, including the finite
initial sizes.
\label{tra:tilted-arrays:fourier-completion}
Its positive-square convention means the matching singular exponent is
$2\sup_d r(d)$. Proposition~\ref{e:nonnormal-moment-conversion} gives
chi-square divergence at most $1/5$ after every integer
$w\ge\sup_d r(d)$ of forward sweeps, uniformly over initial decks.
Since each sweep is $d$ physical shuffles, the stated total-variation
bound follows.
\end{proof}
